% GENERATED FILE. Edit canonical lessons, then run npm run generate:books.
% canonical-source-hash: fnv1a64:4d16763db6f343f4
% canonical-lessons: JA-C418-shomei, JA-C418-youshi, JA-C418-kinyuu, JA-C418-ran, JA-C418-hinichi

\chapter{Filling In a Form}
\label{ch:ja-filling-in-a-form}
\hlchaptermodality{418}

\begin{chapteropening}
\textbf{By the end of this chapter:} \emph{I can fill in a form at a counter.}

Complete the last lesson of chapter 418: I can fill in a form at a counter.
\end{chapteropening}

\section[\texorpdfstring{\ja{しょめい} (shomei)}{しょめい (shomei)}]{\ja{しょめい} (shomei) — a signature}

\label{lesson:JA-C418-shomei}



\begin{quote}
Before the new one: say the Japanese for to carry in one's arms, then the Japanese for to wheel along.
\end{quote}

\subsection*{You'll want to know: \ja{しょめい}}
\textbf{\ja{しょめい}} — \emph{shomei} — \textquotedblleft{}a signature\textquotedblright{}.

\begin{quote}
\textbf{\ja{ここに} \ja{しょめいを} \ja{おねがいします。}} — \emph{koko ni shomei o onegai shimasu} — \textquotedblleft{}Your signature here, please.\textquotedblright{}
\end{quote}

\begin{cousinweb}
\textbf{\ja{しょめい}} (\emph{shomei}) is Sino-Japanese (kango): two Chinese characters read with their on'yomi, readings borrowed from Chinese. They are \emph{sho}, \textquotedblleft{}signature\textquotedblright{} (Middle Chinese \emph{dzyoH}) and \emph{mei}, \textquotedblleft{}name\textquotedblright{} (\emph{mjieng}) (KANJIDIC2; Middle Chinese in Baxter's transcription).
\end{cousinweb}

\begin{grammarlens}[title={What you've built}]
One more word for getting around.
\end{grammarlens}

\subsection*{Your turn}
\begin{itemize}
\raggedright
  \item \emph{Say it:} \emph{shomei}
  \item \emph{Say it:} \emph{shomei}, once more
  \item \emph{Say it:} \emph{koko ni shomei o onegai shimasu} — Your signature here, please.
  \item \emph{Recall:} say the Japanese for to carry in one's arms, then the Japanese for to wheel along, then say \emph{shomei} again
\end{itemize}

\begin{tcolorbox}[breakable,colback=teal!4,colframe=teal!35!black,title={Before you move on}]
What does \ja{しょめい} mean? (\textquotedblleft{}A signature\textquotedblright{}.) Say it once more.
\end{tcolorbox}
\section[\texorpdfstring{\ja{ようし} (yōshi)}{ようし (yōshi)}]{\ja{ようし} (yōshi) — a blank form}

\label{lesson:JA-C418-youshi}



\begin{quote}
Before the new one: say the Japanese for to wheel along, then the Japanese for a signature.
\end{quote}

\subsection*{You'll want to know: \ja{ようし}}
\textbf{\ja{ようし}} — \emph{yōshi} — \textquotedblleft{}a blank form\textquotedblright{}.

\begin{quote}
\textbf{\ja{まどぐちで} \ja{ようしを} \ja{うけとりました。}} — \emph{madoguchi de yōshi o uketorimashita} — \textquotedblleft{}I got a form at the counter.\textquotedblright{}
\end{quote}

\begin{cousinweb}
\textbf{\ja{ようし}} (\emph{yōshi}) is Sino-Japanese (kango): two Chinese characters read with their on'yomi, readings borrowed from Chinese. They are \emph{yō}, \textquotedblleft{}use\textquotedblright{} (Middle Chinese \emph{yowngH}) and \emph{shi}, \textquotedblleft{}paper\textquotedblright{} (\emph{tsyeX}) (KANJIDIC2; Middle Chinese in Baxter's transcription).
\end{cousinweb}

\begin{grammarlens}[title={What you've built}]
One more word for getting around.
\end{grammarlens}

\subsection*{Your turn}
\begin{itemize}
\raggedright
  \item \emph{Say it:} \emph{yōshi}
  \item \emph{Say it:} \emph{yōshi}, once more
  \item \emph{Say it:} \emph{madoguchi de yōshi o uketorimashita} — I got a form at the counter.
  \item \emph{Recall:} say the Japanese for to wheel along, then the Japanese for a signature, then say \emph{yōshi} again
\end{itemize}

\begin{tcolorbox}[breakable,colback=teal!4,colframe=teal!35!black,title={Before you move on}]
What does \ja{ようし} mean? (\textquotedblleft{}A blank form\textquotedblright{}.) Say it once more.
\end{tcolorbox}
\section[\texorpdfstring{\ja{きにゅう} (kinyū)}{きにゅう (kinyū)}]{\ja{きにゅう} (kinyū) — filling in a form; an entry}

\label{lesson:JA-C418-kinyuu}



\begin{quote}
Before the new one: say the Japanese for a signature, then the Japanese for a blank form.
\end{quote}

\subsection*{You'll want to know: \ja{きにゅう}}
\textbf{\ja{きにゅう}} — \emph{kinyū} — \textquotedblleft{}filling in a form; an entry\textquotedblright{}.

\begin{quote}
\textbf{\ja{ようしに} \ja{じゅうしょを} \ja{きにゅうして} \ja{ください。}} — \emph{yōshi ni jūsho o kinyū shite kudasai} — \textquotedblleft{}Please fill in your address on the form.\textquotedblright{}
\end{quote}

\begin{cousinweb}
\textbf{\ja{きにゅう}} (\emph{kinyū}) is Sino-Japanese (kango): two Chinese characters read with their on'yomi, readings borrowed from Chinese. They are \emph{ki}, \textquotedblleft{}account\textquotedblright{} (Middle Chinese \emph{kiH}) and \emph{nyū}, \textquotedblleft{}enter\textquotedblright{} (\emph{nyip}) (KANJIDIC2; Middle Chinese in Baxter's transcription).
\end{cousinweb}

\begin{grammarlens}[title={What you've built}]
One more word for getting around.
\end{grammarlens}

\subsection*{Your turn}
\begin{itemize}
\raggedright
  \item \emph{Say it:} \emph{kinyū}
  \item \emph{Say it:} \emph{kinyū}, once more
  \item \emph{Say it:} \emph{yōshi ni jūsho o kinyū shite kudasai} — Please fill in your address on the form.
  \item \emph{Recall:} say the Japanese for a signature, then the Japanese for a blank form, then say \emph{kinyū} again
\end{itemize}

\begin{tcolorbox}[breakable,colback=teal!4,colframe=teal!35!black,title={Before you move on}]
What does \ja{きにゅう} mean? (\textquotedblleft{}Filling in a form; an entry\textquotedblright{}.) Say it once more.
\end{tcolorbox}
\section[\texorpdfstring{\ja{らん} (ran)}{らん (ran)}]{\ja{らん} (ran) — a box on a form; a column}

\label{lesson:JA-C418-ran}



\begin{quote}
Before the new one: say the Japanese for a blank form, then the Japanese for filling in a form.
\end{quote}

\subsection*{You'll want to know: \ja{らん}}
\textbf{\ja{らん}} — \emph{ran} — \textquotedblleft{}a box on a form; a column\textquotedblright{}.

\begin{quote}
\textbf{\ja{この} \ja{らんに} \ja{なまえを} \ja{かいて} \ja{ください。}} — \emph{kono ran ni namae o kaite kudasai} — \textquotedblleft{}Please put your name in this box.\textquotedblright{}
\end{quote}

\begin{cousinweb}
\textbf{\ja{らん}} (\emph{ran}) is Sino-Japanese (kango): one Chinese character read with its on'yomi, a reading borrowed from Chinese: \emph{ran}, \textquotedblleft{}column\textquotedblright{} (Middle Chinese \emph{lan}) (KANJIDIC2; Middle Chinese in Baxter's transcription).
\end{cousinweb}

\begin{grammarlens}[title={What you've built}]
One more word for getting around.
\end{grammarlens}

\subsection*{Your turn}
\begin{itemize}
\raggedright
  \item \emph{Say it:} \emph{ran}
  \item \emph{Say it:} \emph{ran}, once more
  \item \emph{Say it:} \emph{kono ran ni namae o kaite kudasai} — Please put your name in this box.
  \item \emph{Recall:} say the Japanese for a blank form, then the Japanese for filling in a form, then say \emph{ran} again
\end{itemize}

\begin{tcolorbox}[breakable,colback=teal!4,colframe=teal!35!black,title={Before you move on}]
What does \ja{らん} mean? (\textquotedblleft{}A box on a form; a column\textquotedblright{}.) Say it once more.
\end{tcolorbox}
\section[\texorpdfstring{\ja{ひにち} (hinichi)}{ひにち (hinichi)}]{\ja{ひにち} (hinichi) — a date (for something planned)}

\label{lesson:JA-C418-hinichi}



\begin{quote}
Before the new one: say the Japanese for filling in a form, then the Japanese for a box on a form.
\end{quote}

\subsection*{You'll want to know: \ja{ひにち}}
\textbf{\ja{ひにち}} — \emph{hinichi} — \textquotedblleft{}a date (for something planned)\textquotedblright{}.

\begin{quote}
\textbf{\ja{よやくの} \ja{ひにちを} \ja{きめました。}} — \emph{yoyaku no hinichi o kimemashita} — \textquotedblleft{}We decided on the date of the booking.\textquotedblright{}
\end{quote}

\begin{cousinweb}
\textbf{\ja{ひにち}} (\emph{hinichi}) is built from both layers of the vocabulary: the native Japanese (wago) \emph{hi}, \textquotedblleft{}day\textquotedblright{}, inherited rather than borrowed, and the Sino-Japanese \emph{nichi}, \textquotedblleft{}day\textquotedblright{} (Middle Chinese \emph{nyit}), an on'yomi borrowed from Chinese (JMdict; KANJIDIC2; Middle Chinese in Baxter's transcription).
\end{cousinweb}

\begin{grammarlens}[title={What you've built}]
One more word for getting around.
\end{grammarlens}

\subsection*{Your turn}
\begin{itemize}
\raggedright
  \item \emph{Say it:} \emph{hinichi}
  \item \emph{Say it:} \emph{hinichi}, once more
  \item \emph{Say it:} \emph{yoyaku no hinichi o kimemashita} — We decided on the date of the booking.
  \item \emph{Recall:} say the Japanese for filling in a form, then the Japanese for a box on a form, then say \emph{hinichi} again
\end{itemize}

\begin{tcolorbox}[breakable,colback=teal!4,colframe=teal!35!black,title={Before you move on}]
What does \ja{ひにち} mean? (\textquotedblleft{}A date (for something planned)\textquotedblright{}.) Say it once more.
\end{tcolorbox}
